% GENERATED FILE. Edit canonical lessons, then run npm run generate:books.
% canonical-source-hash: fnv1a64:baf8e4dbf63c1a99
% canonical-lessons: TA-C435-kataloram, TA-C435-valaikuta, TA-C435-munai, TA-C435-uppalam, TA-C435-kulankal

\chapter{Bays and Capes}
\label{ch:ta-bays-and-capes}
\hlchaptermodality{435}

\begin{chapteropening}
\textbf{By the end of this chapter:} \emph{I can talk about the coast, bays, capes, salt pans and pebbles.}

Complete the last lesson of chapter 435: I can talk about the coast, bays, capes, salt pans and pebbles.
\end{chapteropening}

\section[\texorpdfstring{\ta{கடலோரம்} (kaṭalōram)}{கடலோரம் (kaṭalōram)}]{\ta{கடலோரம்} (kaṭalōram) — the coast; the seaside}

\label{lesson:TA-C435-kataloram}



\begin{quote}
Before the new one: say the Tamil for a palmyra palm, then the Tamil for a bush.
\end{quote}

\subsection*{You'll want to know: \ta{கடலோரம்}}
\textbf{\ta{கடலோரம்}} — \emph{kaṭalōram} — \textquotedblleft{}the coast; the seaside\textquotedblright{}.

\begin{quote}
\ta{கடலோரக்} \ta{கிராமங்களில்} \ta{மீனவர்கள்} \ta{வாழ்கிறார்கள்}.
\end{quote}

\begin{quote}
\emph{kaṭalōrak kirāmaṅkaḷil mīṉavarkaḷ vāḻkiṟārkaḷ.}
\end{quote}

\begin{quote}
Fishermen live in the coastal villages.
\end{quote}

As a modifier it drops -\ta{ம்}: \ta{கடலோரக்} \ta{கிராமம்}, a coastal village.

\begin{cousinweb}
\textbf{\ta{கடலோரம்}} (\emph{kaṭalōram}) is built from \ta{கடல்}, \textquotedblleft{}the sea\textquotedblright{} (DEDR 1118), and \ta{ஓரம்}, \textquotedblleft{}an edge, a border\textquotedblright{} (DEDR 1060).
\end{cousinweb}

\begin{grammarlens}[title={What you've built}]
One more word for getting around.
\end{grammarlens}

\subsection*{Your turn}
\begin{itemize}
\raggedright
  \item \emph{Say it:} \emph{kaṭalōram}
  \item \emph{Say it:} \emph{kaṭalōram}, once more
  \item \emph{Say it:} \emph{kaṭalōrak kirāmaṅkaḷil mīṉavarkaḷ vāḻkiṟārkaḷ.} — Fishermen live in the coastal villages.
  \item \emph{Recall:} say the Tamil for a palmyra palm, then the Tamil for a bush, then say \emph{kaṭalōram} again
\end{itemize}

\begin{tcolorbox}[breakable,colback=teal!4,colframe=teal!35!black,title={Before you move on}]
What does \ta{கடலோரம்} mean? (\textquotedblleft{}The coast; the seaside\textquotedblright{}.) Say it once more.
\end{tcolorbox}
\section[\texorpdfstring{\ta{வளைகுடா} (vaḷaikuṭā)}{வளைகுடா (vaḷaikuṭā)}]{\ta{வளைகுடா} (vaḷaikuṭā) — a bay; a gulf}

\label{lesson:TA-C435-valaikuta}



\begin{quote}
Before the new one: say the Tamil for a bush, then the Tamil for the coast.
\end{quote}

\subsection*{You'll want to know: \ta{வளைகுடா}}
\textbf{\ta{வளைகுடா}} — \emph{vaḷaikuṭā} — \textquotedblleft{}a bay; a gulf\textquotedblright{}.

\begin{quote}
\ta{வளைகுடாவில்} \ta{அலைகள்} \ta{அமைதியாக} \ta{இருந்தன}.
\end{quote}

\begin{quote}
\emph{vaḷaikuṭāvil alaikaḷ amaitiyāka iruntaṉa.}
\end{quote}

\begin{quote}
The waves in the bay were calm.
\end{quote}

\ta{வளை}, to bend, + \ta{குடா}, a curve: where the coast bends in.

\begin{cousinweb}
\textbf{\ta{வளைகுடா}} (\emph{vaḷaikuṭā}) is built from \ta{வளை}, \textquotedblleft{}to bend\textquotedblright{} (DEDR 5314), and \ta{குடா}, \textquotedblleft{}a bend, a curve\textquotedblright{} (DEDR 2054).
\end{cousinweb}

\begin{grammarlens}[title={What you've built}]
One more word for getting around.
\end{grammarlens}

\subsection*{Your turn}
\begin{itemize}
\raggedright
  \item \emph{Say it:} \emph{vaḷaikuṭā}
  \item \emph{Say it:} \emph{vaḷaikuṭā}, once more
  \item \emph{Say it:} \emph{vaḷaikuṭāvil alaikaḷ amaitiyāka iruntaṉa.} — The waves in the bay were calm.
  \item \emph{Recall:} say the Tamil for a bush, then the Tamil for the coast, then say \emph{vaḷaikuṭā} again
\end{itemize}

\begin{tcolorbox}[breakable,colback=teal!4,colframe=teal!35!black,title={Before you move on}]
What does \ta{வளைகுடா} mean? (\textquotedblleft{}A bay; a gulf\textquotedblright{}.) Say it once more.
\end{tcolorbox}
\section[\texorpdfstring{\ta{முனை} (muṉai)}{முனை (muṉai)}]{\ta{முனை} (muṉai) — a cape; a point, a tip}

\label{lesson:TA-C435-munai}



\begin{quote}
Before the new one: say the Tamil for the coast, then the Tamil for a bay.
\end{quote}

\subsection*{You'll want to know: \ta{முனை}}
\textbf{\ta{முனை}} — \emph{muṉai} — \textquotedblleft{}a cape; a point, a tip\textquotedblright{}.

\begin{quote}
\ta{நாட்டின்} \ta{தெற்கு} \ta{முனையில்} \ta{மூன்று} \ta{கடல்கள்} \ta{சேர்கின்றன}.
\end{quote}

\begin{quote}
\emph{nāṭṭiṉ teṟku muṉaiyil mūṉṟu kaṭalkaḷ cērkiṉṟaṉa.}
\end{quote}

\begin{quote}
Three seas meet at the southern tip of the country.
\end{quote}

Also the tip of a pen or a knife: \ta{பேனாவின்} \ta{முனை}.

\begin{cousinweb}
\textbf{\ta{முனை}} (\emph{muṉai}) is native Dravidian (DEDR 5020), \textquotedblleft{}a front, a point, a sharpened end\textquotedblright{}, the family of \ta{முன்}, \textquotedblleft{}front\textquotedblright{}.
\end{cousinweb}

\begin{grammarlens}[title={What you've built}]
One more word for getting around.
\end{grammarlens}

\subsection*{Your turn}
\begin{itemize}
\raggedright
  \item \emph{Say it:} \emph{muṉai}
  \item \emph{Say it:} \emph{muṉai}, once more
  \item \emph{Say it:} \emph{nāṭṭiṉ teṟku muṉaiyil mūṉṟu kaṭalkaḷ cērkiṉṟaṉa.} — Three seas meet at the southern tip of the country.
  \item \emph{Recall:} say the Tamil for the coast, then the Tamil for a bay, then say \emph{muṉai} again
\end{itemize}

\begin{tcolorbox}[breakable,colback=teal!4,colframe=teal!35!black,title={Before you move on}]
What does \ta{முனை} mean? (\textquotedblleft{}A cape; a point, a tip\textquotedblright{}.) Say it once more.
\end{tcolorbox}
\section[\texorpdfstring{\ta{உப்பளம்} (uppaḷam)}{உப்பளம் (uppaḷam)}]{\ta{உப்பளம்} (uppaḷam) — a salt pan}

\label{lesson:TA-C435-uppalam}



\begin{quote}
Before the new one: say the Tamil for a bay, then the Tamil for a cape.
\end{quote}

\subsection*{You'll want to know: \ta{உப்பளம்}}
\textbf{\ta{உப்பளம்}} — \emph{uppaḷam} — \textquotedblleft{}a salt pan\textquotedblright{}.

\begin{quote}
\ta{கடலோரத்தில்} \ta{நிறைய} \ta{உப்பளங்கள்} \ta{இருக்கின்றன}.
\end{quote}

\begin{quote}
\emph{kaṭalōrattil niṟaiya uppaḷaṅkaḷ irukkiṉṟaṉa.}
\end{quote}

\begin{quote}
There are many salt pans along the coast.
\end{quote}

Plural: \ta{உப்பளங்கள்}, \emph{uppaḷaṅkaḷ}.

\ta{உப்பு}, salt, + \ta{அளம்}, a salt pan; plural \ta{உப்பளங்கள்}.

\begin{cousinweb}
\textbf{\ta{உப்பளம்}} (\emph{uppaḷam}) is native Dravidian: DEDR 2674 lists it, \textquotedblleft{}saltpans\textquotedblright{}, built from \ta{உப்பு}, \textquotedblleft{}salt\textquotedblright{}, and \ta{அளம்}, \textquotedblleft{}a saltpan\textquotedblright{}.
\end{cousinweb}

\begin{grammarlens}[title={What you've built}]
One more word for getting around.
\end{grammarlens}

\subsection*{Your turn}
\begin{itemize}
\raggedright
  \item \emph{Say it:} \emph{uppaḷam}
  \item \emph{Say it:} \emph{uppaḷam}, once more
  \item \emph{Say it:} \emph{kaṭalōrattil niṟaiya uppaḷaṅkaḷ irukkiṉṟaṉa.} — There are many salt pans along the coast.
  \item \emph{Recall:} say the Tamil for a bay, then the Tamil for a cape, then say \emph{uppaḷam} again
\end{itemize}

\begin{tcolorbox}[breakable,colback=teal!4,colframe=teal!35!black,title={Before you move on}]
What does \ta{உப்பளம்} mean? (\textquotedblleft{}A salt pan\textquotedblright{}.) Say it once more.
\end{tcolorbox}
\section[\texorpdfstring{\ta{கூழாங்கல்} (kūḻāṅkal)}{கூழாங்கல் (kūḻāṅkal)}]{\ta{கூழாங்கல்} (kūḻāṅkal) — a pebble}

\label{lesson:TA-C435-kulankal}



\begin{quote}
Before the new one: say the Tamil for a cape, then the Tamil for a salt pan.
\end{quote}

\subsection*{You'll want to know: \ta{கூழாங்கல்}}
\textbf{\ta{கூழாங்கல்}} — \emph{kūḻāṅkal} — \textquotedblleft{}a pebble\textquotedblright{}.

\begin{quote}
\ta{ஆற்றங்கரையில்} \ta{குழந்தைகள்} \ta{கூழாங்கற்களைச்} \ta{சேகரித்தார்கள்}.
\end{quote}

\begin{quote}
\emph{āṟṟaṅkaraiyil kuḻantaikaḷ kūḻāṅkaṟkaḷaic cēkarittārkaḷ.}
\end{quote}

\begin{quote}
On the riverbank the children collected pebbles.
\end{quote}

Plural: \ta{கூழாங்கற்கள்}, \emph{kūḻāṅkaṟkaḷ}.

\ta{கல்} makes its plural as \ta{கற்கள்}: \ta{கூழாங்கற்கள்}.

\begin{cousinweb}
\textbf{\ta{கூழாங்கல்}} (\emph{kūḻāṅkal}) is built with \ta{கல்}, \textquotedblleft{}a stone, a pebble\textquotedblright{} (DEDR 1298); the origin of its first part is not settled.
\end{cousinweb}

\begin{grammarlens}[title={What you've built}]
One more word for getting around.
\end{grammarlens}

\subsection*{Your turn}
\begin{itemize}
\raggedright
  \item \emph{Say it:} \emph{kūḻāṅkal}
  \item \emph{Say it:} \emph{kūḻāṅkal}, once more
  \item \emph{Say it:} \emph{āṟṟaṅkaraiyil kuḻantaikaḷ kūḻāṅkaṟkaḷaic cēkarittārkaḷ.} — On the riverbank the children collected pebbles.
  \item \emph{Recall:} say the Tamil for a cape, then the Tamil for a salt pan, then say \emph{kūḻāṅkal} again
\end{itemize}

\begin{tcolorbox}[breakable,colback=teal!4,colframe=teal!35!black,title={Before you move on}]
What does \ta{கூழாங்கல்} mean? (\textquotedblleft{}A pebble\textquotedblright{}.) Say it once more.
\end{tcolorbox}
